% =============================================================================
% Figure captions for
%   "A Process Is a Ladder: Catalysis Without Identities"
%
% Six panels, each a single PNG with four charts in a row and at least one
% three-dimensional chart, generated by generate_panels.py from the validation
% JSON in ../validation/results/.  Include with e.g.
%
%   \begin{figure*}[!htbp]
%     \centering
%     \includegraphics[width=\textwidth]{panel_01.png}
%     \captionPanelOne
%     \label{fig:panel1}
%   \end{figure*}
% =============================================================================


% -----------------------------------------------------------------------------
\newcommand{\captionPanelOne}{%
\caption{\textbf{The floor, the hypothesis it requires, and alignment as a
minimum cut.}
\textbf{(A)}~Floor $\beta$ computed---not assumed---as the minimum edge weight
over four finite contact graphs spanning twelve orders of magnitude, including
an adversarial instance whose smallest weight is $10^{-12}$. Every instance
returns a strictly positive floor. Positivity is a property of finiteness and
strict positivity of the weights, not of their scale.
\textbf{(B)}~The failure branch. A system permitted to refine without limit
generates thicknesses $1/n$ (orange): each is strictly positive and the infimum
is zero, with every $\varepsilon$ in a ladder down to $10^{-9}$ eventually
cleared. Truncating the same sequence to a finite system (blue, first ten
stages) restores a positive floor at $0.1$ (dashed). The uniform floor therefore
requires finiteness, and this chart is included so that the requirement is
visible rather than tacit.
\textbf{(C)}~Alignment computed as a minimum residue over admissible cuts,
against an independent brute-force enumeration of every bipartition, on forty
randomly generated graphs. Points lie on the identity to machine precision. This
is a check of the implementation against ground truth, not a restatement of the
definition.
\textbf{(D)}~Floor surface over edge count and weight scale. The surface is
positive everywhere on the domain and descends toward zero only as the scale
does; increasing the edge count at fixed scale lowers the floor but never
abolishes it, which is the content of the minimum over a finite set.}}


% -----------------------------------------------------------------------------
\newcommand{\captionPanelTwo}{%
\caption{\textbf{Probing is free, committing is floored, and a structure
committed once is thereafter traversed at no cost.}
\textbf{(A)}~Residue deposited by a commitment against the measure of the region
committed, over six decades of measure. The deposit is constant and sits at the
floor (dashed): the cost of individuating a region does not depend on how
improbable that region is. This is the first clause of the asymmetry and it is
the clause that makes rarity irrelevant to cost.
\textbf{(B)}~The commitment counter $M$ under the two operations. Commits
increment it one for one; probes leave it flat at zero however many are
performed. The counter is the logical clock of the semantics, and the two lines
are the operational content of the probe/commit distinction.
\textbf{(C)}~Reuse. A single commitment is performed, after which two thousand
traversals are executed. $M$ rises once and then stays; cumulative residue rises
once and then stays. The structure is not consumed by the traversals, and no
energetic argument enters the computation---only the asymmetry.
\textbf{(D)}~Probe frequency against region measure, with commitment cost as the
vertical coordinate and the constant-cost plane overlaid. Frequency tracks
measure over four decades while cost does not vary at all: rarity affects how
often a region is visited and not what visiting it costs.}}


% -----------------------------------------------------------------------------
\newcommand{\captionPanelThree}{%
\caption{\textbf{Composition is multiplicative, and the alternatives are
excluded by comparison rather than by assertion.}
\textbf{(A)}~Predicted against simulated composite power for four candidate
composition laws on five hundred randomly generated ladders. The multiplicative
law lies on the identity; the additive law saturates and overshoots, the
maximum law systematically underestimates, and the mean law fails badly. The
comparison is the test: scoring only the favoured law would not have
distinguished it from the others.
\textbf{(B)}~Mean absolute error by law over four thousand ladders, on a
logarithmic axis. The multiplicative law is exact to machine precision; the
alternatives err by $0.047$, $0.179$ and $0.463$ respectively. The gap spans
sixteen orders of magnitude and is not a matter of fit quality.
\textbf{(C)}~Residual fraction against ladder length for four rung powers, on a
logarithmic axis. Each trace is a straight line, which is the graphical content
of the residual being a product $\prod_i(1-\pi_i)$ rather than a sum.
\textbf{(D)}~Composite power over two rung powers. The surface is
$1-(1-\pi_1)(1-\pi_2)$: it rises to unity along either edge and is symmetric
under exchange, so order does not matter, and it is strictly below the additive
prediction everywhere in the interior.}}


% -----------------------------------------------------------------------------
\newcommand{\captionPanelFour}{%
\caption{\textbf{Repetition has diminishing returns; the dichotomy between
reaching and plateauing is a property of the tail.}
\textbf{(A)}~Composite power under repetition of a single rung, for four rung
powers. Each curve rises strictly, tends to unity (dashed), and reaches it at no
finite $n$. A ladder of near-duplicate rungs therefore cannot attain a demanding
target however long it is made.
\textbf{(B)}~Marginal contribution of the $n$-th repetition, $\pi(1-\pi)^{n-1}$,
on a logarithmic axis. The decline is geometric. The stronger the rung, the
larger the first contribution and the faster the decay---so repetition is least
useful precisely where each copy is individually most effective.
\textbf{(C)}~Residual gap against rungs applied for three power series, indexed
from $i=2$ so that no rung is absolute and the dichotomy is exhibited by tail
behaviour rather than by a degenerate first step. The divergent series
$\sum 1/i$ drives the gap steadily toward zero; the convergent series
$\sum 1/i^2$ and $\sum 2^{-i}$ plateau at finite residual and do not improve
thereafter.
\textbf{(D)}~Rungs required to reach a target, over target and rung power. The
surface diverges as the target approaches unity and as the rung power approaches
zero, and is low only in the corner where strong rungs meet modest targets. This
is the cost model: what a target costs is set jointly by how demanding it is and
by how much each step accomplishes.}}


% -----------------------------------------------------------------------------
\newcommand{\captionPanelFive}{%
\caption{\textbf{Control lies at the strongest rung, not the weakest.}
\textbf{(A)}~Rung power and sensitivity for the four-rung ladder of the worked
demonstration. The two orderings agree: rung~3, of highest power $0.55$, also
carries the highest sensitivity $0.308$. The naive expectation---that the
weakest rung is the bottleneck and therefore the one worth improving---predicts
the opposite ordering and is wrong.
\textbf{(B)}~Sensitivity normalised by the ladder's residual fraction, against
rung power, for nine hundred random ladders. Every point lies on the curve
$1/(1-\pi_j)$, which is exact: sensitivity is $P/(1-\pi_j)$ with $P$ the
residual, so the normalised quantity depends on that rung alone and increases
strictly in its power. The normalisation is necessary because raw sensitivity
also depends on the other rungs, which obscures the relationship across ladders
of differing composition.
\textbf{(C)}~Gain in composite power from a fixed improvement applied to each
rung in turn, at three improvement sizes. Rung~3 returns the most at every size
and rung~4---the weakest---returns the least. Improving the bottleneck is the
worse investment, by roughly a factor of two at $\delta = 0.1$.
\textbf{(D)}~Sensitivity to the third rung of a three-rung ladder, over the
powers of the other two. The surface is $(1-\pi_1)(1-\pi_2)$: sensitivity is
transmission through the remaining rungs, and it is largest when those rungs are
weak. What a marginal gain buys is set by what the rest of the ladder does with
it, which is why the quantity being maximised is transmission and not slack.}}


% -----------------------------------------------------------------------------
\newcommand{\captionPanelSix}{%
\caption{\textbf{Labels move nothing, powers move everything, and the static
verdict agrees with execution.}
\textbf{(A)}~Composite power before and after relabelling the rungs, on four
hundred random ladders. All points lie on the identity. The relabelling is
performed on annotations external to the formalism, since a rung carries exactly
one datum and has no identity field: a test that compared objects carrying names
would be testing the names.
\textbf{(B)}~The control. The same observable under perturbation of a rung
\emph{power} rather than a label. The distribution of changes is entirely
positive and well away from zero. Without this chart panel~A would be consistent
with a constant observable and would carry no information; with it, panel~A is a
statement about labels specifically.
\textbf{(C)}~Static reachability against execution on five thousand generated
programs. Four thousand and seventeen were accepted statically and nine hundred
and eighty-three rejected; in no case did the static verdict disagree with the
executed outcome, marked by the cross at zero. Both verdicts occur in quantity,
so the test could have failed in either direction---a suite in which every
program was accepted would not have tested rejection.
\textbf{(D)}~The worked ladder in three dimensions: rung powers, sensitivities,
and the gap trajectory that descends from unity to the residual $0.1386$ across
the four commitments. Every quantity plotted is computed from the four declared
powers, and nothing about what the rungs are enters the computation or could be
recovered from it.}}
